\begin{mainquestion}{commongateapp}{Common Gate Application}

\questionpoints
In electronics, we occasionally need to read a signal across a long wire (a ``bitline") which tends to have high parasitic capacitance (which we'll model as $C_{BL}$). Assume our bitline has a parasitic capacitance of $C_{BL} = 10\text{ pF}$. The signal source (in this case, a memory cell) acts as a switch with an internal resistance $R_{cell} = 5\text{ k}\Omega$ when active. To read the data, we need to convert the cell's current into a voltage swing using a load resistor $R_L$.

% =========================================================
\begin{partquestion}{rlpullup}
{First, consider a circuit where the bitline is pulled up directly by $R_L = 20\text{ k}\Omega$. Calculate the time constants for both the \textbf{turn-on} phase (active discharge, when the cell switch closes) and the \textbf{turn-off} phase (recovery, when the cell switch opens).

\begin{center}
\begin{circuitikz}[american]
    \draw (0,0) node[vcc] {$V_{DD}$} to[R, l=$R_L$] (0,-2);
    \draw (0,-2) to[short, -*] (1.5,-2) node[right] {Bitline};
    \draw (0,-2) to[C, l_=$C_{BL}$] (0,-4) node[ground] {};
    \draw (0,-2) -- (-2.5,-2) to[nos, l=Cell Switch] (-2.5,-3) to[R, l=$R_{cell}$] (-2.5,-4) node[ground] {};
\end{circuitikz}
\end{center}}

\begin{answer}[2.5in]
    \textbf{Turn-on (Active Discharge):} 
    When the cell switch closes, the total resistance seen by the capacitor $C_{BL}$ is the parallel combination of the load and the cell. 
    \[\tau_{\text{on}} = (R_L \parallel R_{cell}) \cdot C_{BL}\]
    \[\tau_{\text{on}} = \left(\frac{20\text{ k}\Omega \cdot 5\text{ k}\Omega}{20\text{ k}\Omega + 5\text{ k}\Omega}\right) \cdot 10\text{ pF} = \left(\frac{100}{25}\right)\text{ k}\Omega \cdot 10\text{ pF} = 4\text{ k}\Omega \cdot 10\text{ pF} = 40\text{ ns}\]
    
    \textbf{Turn-off (Recovery):}
    When the cell switch opens, the cell is disconnected. The bitline must now recover to $V_{DD}$ through $R_L$ alone. 
    \[\tau_{\text{off}} = R_L \cdot C_{BL} = 20\text{ k}\Omega \cdot 10\text{ pF} = 200\text{ ns}\]
\end{answer}

\pagebreak
\end{partquestion}

% =========================================================
\begin{partquestion}{commongate}
{To speed up the read, we insert an NMOS transistor in a Common Gate configuration, and add a current source to the bitline to bias the transistor. Let's assume that the current source is $250\mu A$ and the $V_{OV}$ of the transistor is 100mV. The heavily-capacitive bitline is connected to the \textbf{Source}, and $R_L$ is moved to the \textbf{Drain}. 

\begin{center}
\begin{circuitikz}[american, arrowmos]
    \draw (0,0) node[nmos] (mos) {};
    \draw (mos.G) node[vcc]{$V_{bias}$} {};
    \draw (mos.D) to[R, l=$R_L$] (0, 2) node[vcc] {$V_{DD}$};
    \draw (mos.D) to[short, -o] (1, 0.77) node[right] {$V_{out}$};
    
    % Connect Source to the Bitline node
    \draw (mos.S) -- (0,-1.5) coordinate (BL);
    \draw (BL) to[short, *-] (1.5,-1.5) node[right] {Bitline};
    \draw (1,-1.5) to [I]  (1,-4) node[ground]{};
    \draw (BL) to[C, l_=$C_{BL}$] (0,-4) node[ground] {};
    
    % Add the memory cell switch and resistance (Identical to Part a)
    \draw (BL) -- (-2.5,-1.5) to[nos, l=Cell Switch] (-2.5,-2.5) to[R, l=$R_{cell}$] (-2.5,-4) node[ground] {};
\end{circuitikz}
\end{center}

Before calculating the new speed, derive the input resistance $R_{in}$ looking into the Source of the transistor. Assume there is no body effect ($g_{mb} = 0$) and no channel length modulation ($r_o = \infty$).
\\
\\
\textit{Hint: To find the input resistance, we 'look into' the transistor's Source terminal, temporarily ignoring the external bitline capacitance and memory cell.}}

\begin{answer}[3.5in]
    To find $R_{in}$, we apply a small-signal test voltage $v_x$ to the Source and measure the test current $i_x$ entering the node. 
    
    \begin{center}
    \begin{circuitikz}[american, scale=0.9, transform shape]
        % Test Source
        \draw (0,-2) node[ground] {} to[V, l=$v_x$, invert] (0,0) to[short, i=$i_x$] (1.5,0);
        
        % Gate
        \draw (2,2) to[short, *-] (1,2) node[ground] {};
        \node[above] at (2,2) {G};
        
        % Source
        \draw (1.5,0) to[short, *-*] (4,0);
        \node[below right] at (2,0) {S};
        
        % Vgs
        \draw (2,2) to[open, v=$v_{gs}$] (2,0);
        
        % Drain & Dependent Source
        \draw (4,2) to[cI, l=$g_m v_{gs}$] (4,0);
        \node[above] at (4,2) {D};
        
        % Load
        \draw (4,2) to[short, *-] (6,2) to[R, l=$R_L$] (6,-2) node[ground] {};
    \end{circuitikz}
    \end{center}
    
    The gate is connected to AC ground, so the small-signal gate-to-source voltage is:
    \[v_{gs} = v_g - v_s = 0 - v_x = -v_x\]
    
    The transistor's dependent current source pushes a current of $g_m v_{gs}$ from drain to source. 
    
    Applying KCL at the Source node (sum of currents entering = 0):
    \[i_x + g_m v_{gs} = 0 \implies i_x + g_m(-v_x) = 0 \implies i_x = g_m v_x\]
    
    Therefore, the input resistance is:
    \[R_{in} = \frac{v_x}{i_x} = \frac{1}{g_m}\]
    
    To find the numerical value, we calculate $g_m$ from the bias conditions:
    \[g_m = \frac{2I_D}{V_{OV}} = \frac{2(250\mu\text{A})}{100\text{ mV}} = 5\text{ mS}\]
    \[R_{in} = \frac{1}{5\text{ mS}} = 200\,\Omega\]
\end{answer}

\qnewpage
\end{partquestion}

% =========================================================
\begin{partquestion}{switchingtimes}
{Using your result from Part (b), calculate the new turn-on and turn-off time constants ($\tau_{\text{on}}$ and $\tau_{\text{off}}$) at the bitline node. Assume $R_{in}$ remains constant during both transitions. Briefly compare your results to the Passive Approach.}

\begin{answer}[3.5in]
    \textbf{Turn-on (Active Discharge):} 
    When the cell switch closes, the total resistance discharging the bitline capacitance is the parallel combination of the cell's internal resistance ($R_{cell}$) and the transistor's input resistance ($R_{in}$). 
    \[R_{total,on} = R_{in} \parallel R_{cell}\]
    \[R_{total,on} = 200\,\Omega \parallel 5000\,\Omega = \frac{200 \cdot 5000}{200 + 5000}\,\Omega \approx 192.3\,\Omega\]
    \[\tau_{\text{on}} = R_{total,on} \cdot C_{BL} \approx 192.3\,\Omega \cdot 10\text{ pF} \approx 1.92\text{ ns}\]
    
    \vspace{1em}
    \textbf{Turn-off (Recovery):}
    When the cell switch opens, the memory cell is disconnected. The bitline capacitance now only sees the input impedance of the Source ($R_{in}$).
    \[\tau_{\text{off}} = R_{in} \cdot C_{BL} = 200\,\Omega \cdot 10\text{ pF} = 2\text{ ns}\]
    
    \vspace{1em}
    \textbf{Comparison:}
    In the passive approach, the turn-off phase created a massive 200 ns bottleneck. By using the Common Gate configuration, the turn-on speed increases by roughly 20x (40 ns $\rightarrow$ 1.92 ns), and the turn-off bottleneck is completely eliminated, yielding a 100x speedup (200 ns $\rightarrow$ 2 ns).
\end{answer}

\end{partquestion}

% =========================================================
\begin{partquestion}{timeconstant}
{Show that in this new topology, the time constant now doesn't depend on the load resistance. Given that delay doesn't limit the load resistance, what does?}

\begin{answer}[3in]
    As derived in Part (b) and (c), the time constant at the bitline node is determined by the capacitance $C_{BL}$ and the equivalent resistance seen at the source node $(1/g_m \parallel R_{cell})$. 
    
    Because the load resistor $R_L$ is located at the drain, and we assume an ideal transistor ($r_o = \infty$), $R_L$ is isolated from the input node. It does not appear anywhere in the RC time constant equation for the bitline, making the delay independent of the load resistance. 
    
    While delay no longer limits the maximum size of $R_L$, \textbf{voltage headroom} does. If $R_L$ is made too large, the DC voltage drop across it will become too great, lowering the drain voltage and eventually pushing the transistor out of saturation (into the triode region).
\end{answer}

\pagebreak
\end{partquestion}

% =========================================================
\begin{partquestion}{increaserl}
{Let's increase $R_L$ from $20\text{ k}\Omega$ to $100\text{ k}\Omega$. Calculate the new active discharge (turn-on) time constant for both the \textbf{Passive Approach} and the \textbf{Common Gate Approach}. Based on your calculations, analyze the results. What does this mathematically demonstrate about the Common Gate topology?}

\begin{answer}[3in]
    \textbf{Passive Approach Calculation:}
    \[\tau_{\text{on}} = (100\text{ k}\Omega \parallel 5\text{ k}\Omega) \cdot 10\text{ pF}\]
    \[\tau_{\text{on}} = \left(\frac{100 \cdot 5}{105}\right)\text{ k}\Omega \cdot 10\text{ pF} \approx 4.76\text{ k}\Omega \cdot 10\text{ pF} = 47.6\text{ ns}\]
    \textit{(Note: The turn-off recovery time dramatically worsens to $100\text{ k}\Omega \cdot 10\text{ pF} = 1\,\mu\text{s}$).}
    
    \textbf{Common Gate Approach Calculation:}
    \[\tau_{\text{on}} = R_{total,on} \cdot C_{BL} = (R_{in} \parallel R_{cell}) \cdot C_{BL}\]
    \[\tau_{\text{on}} \approx 192.3\,\Omega \cdot 10\text{ pF} \approx 1.92\text{ ns}\]
    
    \textbf{Analysis / Conclusion:} 
    In the passive topology, attempting to increase the voltage swing by increasing $R_L$ directly increases the RC time constant, severely slowing down the circuit. 
    
    In the Common Gate topology, the heavily-capacitive bitline node is isolated at the Source, while the high-resistance load is located at the Drain. Because $\tau_{\text{on}}$ relies purely on the parallel combination of $1/g_m$ and $R_{cell}$ discharging $C_{BL}$, changing $R_L$ has zero effect on the input time constant. This mathematically demonstrates that the Common Gate amplifier successfully \textbf{decouples} the massive input capacitance from the load resistor, allowing us to increase the output swing without sacrificing speed.
\end{answer}

\end{partquestion}

\end{mainquestion}
